\section{Conclusion and analytical status}
\label{sec:status-tests}
% Status: derived/open. Summarizes established results, assumptions, and theorem targets.

The construction begins with an exact two-branch algebra.  For
\[
  Q(J)=\mu^2
  \begin{pmatrix}
  0&\sqrt J\\
  \sqrt J&-J
  \end{pmatrix},
  \qquad
  J=s(s+1),
\]
the dimensionless eigenvalues obey \(x^2+Jx-J=0\).  The quotient
\[
  1-\frac{x_+(3/4)}{x_+(2)}
\]
is therefore a rigid output of the construction once the ordered samples
\((3/4,2)\) have been chosen.  The mathematical statement is derived.  The
physical statement begins with the additional assignment and placement
hypotheses.

The electroweak comparison used in this paper is the complex-pole mass ratio
\[
  \sPole=1-\frac{M_{W,\rm pole}^2}{M_{Z,\rm pole}^2}.
\]
This choice gives the spectral language a low-energy meaning and separates the
proposal from \(\MSbar\) running angles, effective weak angles, Breit--Wigner
line-shape masses, and high-scale GUT boundary values.  The pole placement is
an assumption awaiting a dynamical pole condition.  Appendix~\ref{app:target-pole-placement}
states the corresponding theorem target.

The Standard Model gauge-Higgs sector supplies the projective datum
\[
  \frac{M_W^2}{M_Z^2}=\frac{g^2}{g^2+g'^2}
\]
and the broken-to-unbroken vacuum ray.  It also supplies two representation
data relevant to the proposed assignment: the Higgs/order-parameter doublet
with \(J_H=3/4\), and the adjoint/current channel with \(J_{\rm adj}=2\).
The required calculation is the ordered map
\[
  (J_W,J_Z)=(J_H,J_{\rm adj})=(3/4,2)
\]
from gauge-Higgs, boundary, endpoint, or compactification data.  This map is a
central theorem target of the paper.

The negative branch belongs to the same diagonalization as the positive branch.
Its scalar interpretation remains conditional on a gauge-invariant functional,
a normalization, and a scheme.  Candidate targets include a Higgs-potential
parameter, a scalar pole, a boundary modulus, a brane separation mode, or a
compactification deformation.  This branch is retained because a source theory
that derives the positive-branch determinant must also account for the partner
eigenvalue.

The string, endpoint, interval, and \(G_2\) material gives mechanisms for a
two-channel kernel.  In common notation the required reduction is
\[
  K_J(\lambda)=
  \begin{pmatrix}
  \lambda+\Sigma_{hh,J} & \Sigma_{ha,J}\\
  \Sigma_{ah,J} & \lambda+\Sigma_{aa,J}
  \end{pmatrix},
  \qquad
  \Sigma_{hh,J}=0,\quad
  \Sigma_{aa,J}=J,\quad
  \Sigma_{ha,J}\Sigma_{ah,J}=J .
\]
Open-string endpoint data, interval boundary conditions, and localized
\(G_2\) gauge sectors provide source-backed addresses for such a kernel.  The
exact DeVries entries remain to be derived from one of those source theories.

The paper therefore establishes a precise algebraic clue and a constrained
research program around it.  A complete derivation requires three linked
results: the two-channel determinant from a source theory, the ordered
electroweak assignment, and a controlled scalar interpretation of the negative
branch.  Additional completion data include the pole-scheme audit, the
low-energy placement theorem, and compatibility with the global form of the
Standard Model gauge group.  Appendix~\ref{app:theorem-targets} collects these
requirements in theorem-target form.
